\chapter{Mở đầu}
\label{chap:mo-dau}

\section{Đặt vấn đề}
\label{sec:dat-van-de}

Định vị lỗi là bước tốn công nhất của gỡ lỗi: khi một bộ kiểm thử báo thất bại, lập trình viên phải tìm ra dòng mã gây
lỗi trong hàng trăm dòng có liên quan. Họ phương pháp được dùng rộng rãi nhất là định vị lỗi dựa trên phổ độ phủ
(SBFL)~\cite{jones2002,abreu2009,wong2016survey}: một dòng càng hay được các kiểm thử thất bại chạy qua và càng ít được
các kiểm thử thành công chạy qua thì càng đáng nghi. Cách làm này rẻ, không cần hiểu ngữ nghĩa chương trình, và hoạt
động tốt khi lỗi làm thay đổi đường thực thi.

Với \emph{lỗi ngữ nghĩa} --- chương trình chạy đúng đường nhưng tính sai giá trị, ví dụ một phép so sánh lệch biên hay
một hằng số trừ thừa một đơn vị --- kiểm thử thành công và thất bại thường chạy qua cùng một tập dòng. Phổ độ phủ khi
đó gần như không phân biệt được dòng lỗi với các dòng xung quanh. Trên bộ dữ liệu SemFault-250, SBFL theo công thức
Ochiai chỉ đặt đúng dòng lỗi ở hạng 1 trong \RqOneSbflOchiaiFairAllAccOne\% số chương trình khi được chấm công bằng
(Chương~\ref{chap:ket-qua}).

Mô hình ngôn ngữ lớn (LLM) mở ra hướng ngược lại: đọc mã nguồn, hiểu ý đồ và đoán dòng sai. Các hệ thống định vị lỗi
dựa trên LLM như AutoFL~\cite{kang2024autofl}, AgentFL~\cite{qin2025agentfl} hay FlexFL~\cite{xu2025flexfl} cho kết quả
ấn tượng, nhưng câu trả lời của LLM mang tính ngẫu nhiên, khó kiểm chứng và có thể chịu ảnh hưởng của việc mô hình đã
thấy dữ liệu đánh giá khi huấn luyện~\cite{sallou2024,ramos2024memorize}.

SemLoc~\cite{yang2026semloc} --- bài báo trung tâm trong ba bài được giảng viên giao --- đề xuất một cách dung hòa: bắt
LLM viết ra các \emph{ràng buộc ngữ nghĩa có thể thực thi}, chèn chúng vào chương trình, chạy kiểm thử, rồi dùng phổ
vi phạm ràng buộc để xếp hạng dòng. Ý tưởng cốt lõi là \emph{neo lập luận của LLM vào bằng chứng thực thi}
(structured grounding): LLM đưa giả thuyết, còn kết quả chạy chương trình mới quyết định. SemLoc báo cáo Acc@1
\RqOneSemlocClaudeTZeroEightAuthorAllAccOne\% trên SemFault-250, vượt xa SBFL.

\section{Câu hỏi nghiên cứu}
\label{sec:cau-hoi}

Tiểu luận đặt một câu hỏi đơn giản cho SemLoc: \textbf{độ chính xác của nó đến từ bằng chứng thực thi (``grounded'')
hay từ việc LLM đoán (``guessed'')?} Câu hỏi này được tách thành ba câu hỏi nghiên cứu:

\begin{itemize}
\item \textbf{RQ1 -- Tái hiện và kiểm tra cách đánh giá.} Kết quả công bố của SemLoc có tái hiện được không, và còn
đứng vững khi chấm điểm công bằng, khi bỏ các thành phần không được mô tả trong bài báo, khi sửa các dòng lỗi chuẩn
(ground truth) sai của bộ dữ liệu?
\item \textbf{RQ2 -- Cải tiến và đối chứng.} Bổ sung tín hiệu độ phủ vào SemLoc (cải tiến CovGate) có làm tốt hơn
không, và một prompt duy nhất chứa bằng chứng thực thi --- kiểm thử thất bại và độ phủ (LLM + độ phủ) --- so với toàn bộ quy trình nhiều bước của SemLoc thì sao, về cả
độ chính xác lẫn chi phí?
\item \textbf{RQ3 -- Độ bền.} Các kết luận trên có giữ được trên lỗi thật (BugsInPy~\cite{widyasari2020}), trên mã đã
đổi tên định danh (kiểm tra việc LLM ghi nhớ bộ dữ liệu) và trên năm họ mô hình khác nhau không?
\end{itemize}

\section{Đóng góp chính}
\label{sec:dong-gop-tom-tat}

Tiểu luận có các đóng góp sau; mỗi đóng góp gắn với số liệu được sinh tự động từ mã nguồn công khai.

\begin{enumerate}
\item \textbf{Đo phần ``grounded'' và phần ``guessed''.} Mã nguồn công bố cộng cho dòng LLM tự đoán một điểm thưởng
lớn hơn mọi điểm ràng buộc, nên dòng đầu tiên của SemLoc \emph{luôn} là dòng LLM đoán: Acc@1 của SemLoc
(\RqTwoRepairedAllSemlocAccOneMean\%) chính là của dòng đoán, còn riêng phổ vi phạm ràng buộc chỉ đạt
\RqTwoRepairedAllSemlocNoPriorAccOneMean\%. Bằng chứng thực thi phát huy ở top-3: nó nâng Acc@3 từ
\RqTwoRepairedAllSemlocAccOneMean\% (chỉ dòng đoán) lên \RqTwoRepairedAllSemlocAccThreeMean\%.
\item \textbf{Một prompt có bằng chứng thực thi vượt quy trình nhiều bước trên các lỗi được chạy tới.} Phương pháp đối
chứng LLM + độ phủ do nhóm thiết kế (một lượt gọi, có tên kiểm thử thất bại, thông báo lỗi và độ phủ từng dòng) hơn
SemLoc \RqTwoOriginalExecutableLlmCoverageDAccThree{} điểm Acc@3 trên các lỗi mà kiểm thử thất bại chạy tới (bộ dòng lỗi
chuẩn gốc), tốn ít token hơn khoảng năm lần và không thay đổi đáng kể khi đổi tên định danh; trên toàn bộ nhãn gốc, gồm
cả lỗi không kiểm thử nào chạy tới, mức hơn chỉ còn \RqTwoOriginalAllLlmCoverageDAccThree{} điểm và SemLoc tốt hơn ở Acc@5.
\item \textbf{Kiểm tra cách đánh giá.} Tính lại khớp toàn bộ cấu hình từ nhật ký công bố, nhưng chỉ ra điểm thưởng ẩn
cho dòng LLM đoán, bước kiểm chứng phản thực tế không làm thay đổi thứ hạng nào, mã chấm điểm cho gỡ lỗi chia đôi
\RqOneDdAuthorAllAccOne\% trong khi bài báo ghi 0, hai lỗi trong bộ chạy BugsInPy, và chạy lại với chính mô hình của bài
báo không đạt con số trong nhật ký.
\item \textbf{SemFault-250 đã sửa.} Dòng lỗi chuẩn được kiểm chứng lại bằng thực thi: \GtConfirmedN{} được xác nhận,
\GtCorrectedN{} được sửa vị trí, \GtUnverifiedN{} chưa kiểm chứng được và giữ nhãn gốc; \GtBlankN{} dòng chuẩn gốc là
dòng trống.
\item \textbf{CovGate}, một bước phá hòa bằng Ochiai cho SemLoc với mức tăng nhỏ ở top-3, cùng bảng kiểm tra đánh giá
bảy mục áp dụng cho mọi phương pháp.
\end{enumerate}

\IfAtLeast{lan2}{Ngoài ra, phân tích thất bại (Chương~\ref{chap:that-bai}) cho thấy độ chính xác cấp dòng phụ thuộc
mạnh vào việc dòng lỗi chuẩn được gán ở dòng điều kiện hay ở câu lệnh mà điều kiện đó bảo vệ, và Chương~\ref{chap:dong-gop}
trình bày quan điểm khoa học cùng các mối đe dọa đến tính hợp lệ của kết quả.}

\section{Cấu trúc báo cáo}
\label{sec:cau-truc}

Chương~\ref{chap:tong-quan} tổng quan tài liệu và bộ từ khóa, trong đó phân tích ba bài báo được giao.
Chương~\ref{chap:co-so} trình bày cơ sở lý thuyết về SBFL, lỗi ngữ nghĩa, SemLoc và các công cụ thống kê.
Chương~\ref{chap:phuong-phap} mô tả câu hỏi nghiên cứu, dữ liệu, các phương pháp so sánh và thiết kế thực nghiệm.
Chương~\ref{chap:ket-qua} trả lời ba câu hỏi nghiên cứu.
\IfAtLeast{lan2}{Chương~\ref{chap:that-bai} phân tích các ca thất bại; Chương~\ref{chap:dong-gop} tổng hợp đóng góp,
triết lý khoa học và các mối đe dọa.}
Chương~\ref{chap:ket-luan} kết luận và nêu hướng phát triển.
\IfAtLeast{full}{Phụ lục hướng dẫn tái hiện toàn bộ kết quả.}
Mã nguồn, dữ liệu và nhật ký gọi LLM được công bố tại \url{https://github.com/doanan293/grounded-or-guessed}.
